\chapter{Sampling Distributions and Standard Error}
\label{ch:sampling-distributions}

Chapter 5 showed how to draw a sample that is random, so that it is
representative on average. Random has a price: a different random sample gives
a different answer. You measured the mean latency of 64 requests and got
48.7 ms. Run the same measurement tomorrow on 64 different requests and you will
get a different number. How different? That question has a precise answer: the
\textbf{standard error}, the typical size of the gap between your estimate and
the truth. Every confidence interval and test later in the book is built from it.

We need three facts about expectation and variance, one change of viewpoint (a
statistic is a random variable), and one simulation. Chapters 7 to 11 rest on
what follows.

\section{Three tools: expectation and variance rules}
\label{sec:sd-tools}

A \textbf{random variable} $X$ is a number whose value is decided by chance: the
roll of a die, the latency of the next request, the height of a randomly chosen
person. Its \textbf{expectation} $\E[X]$ is the long-run average value, and its
\textbf{variance} $\Var(X)=\E[(X-\E[X])^2]$ is the long-run average squared distance
from that centre. The \textbf{standard deviation} $\SD(X)=\sqrt{\Var(X)}$ puts
that spread back in the original units. Two random variables are
\textbf{independent} if knowing the value of one tells you nothing about the
other, like two separate dice. The proofs of the rules below are in
\emph{Probability From the Ground Up}; here we use them as tools.

\begin{example}[One fair die]
$X\in\{1,\dots,6\}$, each with probability $1/6$.
\[
\E[X]=\frac{1+2+3+4+5+6}{6}=3.5,\qquad
\Var(X)=\frac{(2.5^2+1.5^2+0.5^2)\cdot 2}{6}=\frac{17.5}{6}=\frac{35}{12}\approx 2.917.
\]
\end{example}

\begin{theorem}[The three rules]
\label{thm:sd-rules}
For random variables $X,Y$ and constants $a,b$:
\begin{enumerate}
  \item \textbf{Expectation is linear:} $\E[aX+bY]=a\E[X]+b\E[Y]$. This holds
  \emph{always}, independent or not.
  \item \textbf{Scaling squares the variance:} $\Var(aX+b)=a^2\Var(X)$, so
  $\SD(aX+b)=\abs{a}\SD(X)$. Adding a constant shifts but does not spread.
  \item \textbf{Variances of independent variables add:} if $X$ and $Y$ are
  independent, $\Var(X+Y)=\Var(X)+\Var(Y)$ and also
  $\Var(X-Y)=\Var(X)+\Var(Y)$.
\end{enumerate}
\end{theorem}

Tiny checks. Converting latency from ms to seconds is $Y=X/1000$. If
$\SD(X)=30$ ms then $\SD(Y)=0.03$ s and $\Var(Y)=900/10^6=0.0009$ s$^2$. Adding a
fixed 5 ms network overhead changes the mean by 5 and the SD by nothing.

Now roll two independent dice $X_1,X_2$. Rule 3 gives
$\Var(X_1+X_2)=35/12+35/12=35/6\approx 5.833$. Compare that with doubling a single
roll, $2X_1$. Rule 2 gives $\Var(2X_1)=4\cdot 35/12=35/3\approx 11.667$. Same
mean (7), twice the variance. Figure~\ref{fig:sd-dice} shows why: two independent
dice partly cancel each other's extremes; one die doubled never cancels.

\begin{figure}[H]
\centering
\begin{tikzpicture}
\begin{axis}[at={(0,0)}, width=4.6cm, height=4cm, ybar, bar width=4pt,
  ymin=0, ymax=0.2, ytick={0,0.1,0.2}, yticklabel style={/pgf/number format/fixed}, xmin=0.5, xmax=6.5, xtick={1,...,6},
  title={one die $X$}, ylabel={\small probability}, title style={font=\small},
  tick label style={font=\scriptsize}, axis lines=left]
\addplot[fill=histfill, draw=distblue] coordinates {(1,0.1667)(2,0.1667)(3,0.1667)(4,0.1667)(5,0.1667)(6,0.1667)};
\end{axis}
\begin{axis}[at={(4.3cm,0)}, width=4.6cm, height=4cm, ybar, bar width=3pt,
  ymin=0, ymax=0.2, ytick={0,0.1,0.2}, yticklabel style={/pgf/number format/fixed}, xmin=1, xmax=13, xtick={2,7,12},
  title={$X_1+X_2$: Var $5.83$}, title style={font=\small},
  tick label style={font=\scriptsize}, axis lines=left]
\addplot[fill=targetgreen!60, draw=targetgreen] coordinates {(2,0.0278)(3,0.0556)(4,0.0833)(5,0.1111)(6,0.1389)(7,0.1667)(8,0.1389)(9,0.1111)(10,0.0833)(11,0.0556)(12,0.0278)};
\end{axis}
\begin{axis}[at={(8.6cm,0)}, width=4.6cm, height=4cm, ybar, bar width=3pt,
  ymin=0, ymax=0.2, ytick={0,0.1,0.2}, yticklabel style={/pgf/number format/fixed}, xmin=1, xmax=13, xtick={2,7,12},
  title={$2X_1$: Var $11.67$}, title style={font=\small},
  tick label style={font=\scriptsize}, axis lines=left]
\addplot[fill=trapred!50, draw=trapred] coordinates {(2,0.1667)(4,0.1667)(6,0.1667)(8,0.1667)(10,0.1667)(12,0.1667)};
\end{axis}
\end{tikzpicture}
\caption{Both $X_1+X_2$ and $2X_1$ have mean 7. The sum of two independent dice
piles up in the middle because highs and lows tend to offset. Doubling one die
keeps every extreme, so its variance is twice as large.}
\label{fig:sd-dice}
\end{figure}

\begin{trapbox}
Variances add; standard deviations do not. For two independent dice,
$\SD(X_1+X_2)=\sqrt{5.833}=2.415$, not $1.708+1.708=3.416$. And for
\emph{dependent} variables even variances do not simply add: a covariance term
appears (Chapter 4). Every SE formula in this chapter assumes independent
observations.
\end{trapbox}

\section{A statistic is a random variable}
\label{sec:sd-statistic}

A \textbf{statistic} is any number computed from a sample: the mean, the median,
the proportion of conversions, the p99 latency, the slope of a regression line.
Once you have your data, a statistic is just a number, $\xbar=48.7$. Before you
collect the data, the sample $X_1,\dots,X_n$ is random, so the statistic
computed from it is random too. We write it with capital letters, $\Xbar$, to mark
that.

\begin{definition}[Sampling distribution]
The \textbf{sampling distribution} of a statistic is its probability distribution
over all the random samples of size $n$ that could have been drawn. The
\textbf{standard error} of the statistic is the standard deviation of its sampling
distribution.
\end{definition}

Figure~\ref{fig:sd-flow} shows the thought experiment: draw a sample, compute
the statistic, put the sample back, repeat many times, and histogram the results.

\begin{figure}[H]
\centering
\begin{tikzpicture}[
  box/.style={draw, rounded corners, align=center, font=\small, minimum height=0.8cm},
  arr/.style={-{Stealth}, thick, gray}]
\node[box, fill=histfill!30, minimum width=2.4cm, minimum height=2.6cm] (pop) at (0,0)
  {population\\[2pt] mean $\mu$\\ SD $\sigma$};
\node[box, fill=white] (s1) at (3.6,1.3) {sample 1 ($n$)};
\node[box, fill=white] (s2) at (3.6,0) {sample 2 ($n$)};
\node[box, fill=white] (s3) at (3.6,-1.3) {sample 3 ($n$)};
\node[font=\small] at (3.6,-2.05) {$\vdots$};
\node[box, fill=quotientorange!15] (m1) at (6.5,1.3) {$\xbar_1=52.1$};
\node[box, fill=quotientorange!15] (m2) at (6.5,0) {$\xbar_2=46.8$};
\node[box, fill=quotientorange!15] (m3) at (6.5,-1.3) {$\xbar_3=49.9$};
\node[font=\small] at (6.5,-2.05) {$\vdots$};
\foreach \i in {1,2,3} {\draw[arr] (pop.east) -- (s\i.west); \draw[arr] (s\i) -- (m\i);}
\begin{scope}[shift={(8.4,-1.2)}]
  \foreach \x/\h in {0/0.3,0.4/0.9,0.8/1.8,1.2/2.3,1.6/1.6,2.0/0.8,2.4/0.3}
    \draw[fill=quotientorange!60, draw=white] (\x,0) rectangle (\x+0.4,\h);
  \draw[thick] (0,0) -- (2.8,0);
  \node[font=\small, align=center] at (1.4,-0.45) {histogram of $\xbar$};
\end{scope}
\draw[arr] (m2.east) -- (8.3,0);
\end{tikzpicture}
\caption{The sampling distribution is the histogram you would get from repeating
the whole study many times. In practice you run it once; theory (or
simulation) tells you what the histogram would look like.}
\label{fig:sd-flow}
\end{figure}

\subsection{A sampling distribution you can list by hand}

Take a population of three numbers, $\{2,4,9\}$. Its mean is $\mu=5$ and its
variance (dividing by the population size 3) is
\[
\sigma^2=\frac{(2-5)^2+(4-5)^2+(9-5)^2}{3}=\frac{9+1+16}{3}=\frac{26}{3}\approx 8.667.
\]
Draw a sample of $n=2$ \emph{with replacement}, so the two draws are independent.
There are $3\times 3=9$ equally likely ordered samples:

\begin{center}
\begin{tabular}{lccccccccc}
\toprule
sample & (2,2) & (2,4) & (4,2) & (2,9) & (9,2) & (4,4) & (4,9) & (9,4) & (9,9)\\
$\xbar$ & 2 & 3 & 3 & 5.5 & 5.5 & 4 & 6.5 & 6.5 & 9\\
\bottomrule
\end{tabular}
\end{center}

The average of the nine means is $45/9=5=\mu$. Their variance is
\[
\frac{(-3)^2+2(-2)^2+(-1)^2+2(0.5)^2+2(1.5)^2+4^2}{9}
=\frac{9+8+1+0.5+4.5+16}{9}=\frac{39}{9}\approx 4.333=\frac{\sigma^2}{2}.
\]
The mean of $\Xbar$ equals $\mu$ and its variance is half the population variance.
Section~\ref{sec:sd-derive} shows this is no coincidence.

\begin{recoverybox}
Pause and think: if you drew samples of size $n=3$ from $\{2,4,9\}$ instead, what
would you guess the variance of $\Xbar$ to be? (Answer at the end of
Section~\ref{sec:sd-derive}.)
\end{recoverybox}

\begin{figure}[H]
\centering
\begin{tikzpicture}[x=1.15cm]
\draw[densityred, dashed, thick] (5,0.9) -- (5,-1.6);
\node[densityred, font=\small, anchor=south] at (5,0.9) {$\mu=5$};
\draw[-{Stealth}] (1.3,0) -- (9.8,0);
\foreach \t in {2,...,9} \draw (\t,0.08) -- (\t,-0.08) node[below, font=\scriptsize, fill=white, inner sep=1pt] {\t};
\node[anchor=west, font=\small] at (-0.6,0.45) {population};
\foreach \x in {2,4,9} \fill[distblue] (\x,0.45) circle (3pt);
\begin{scope}[yshift=-1.6cm]
\draw[-{Stealth}] (1.3,0) -- (9.8,0);
\foreach \t in {2,...,9} \draw (\t,0.08) -- (\t,-0.08) node[below, font=\scriptsize] {\t};
\node[anchor=west, font=\small] at (-0.6,0.45) {9 sample means};
\foreach \x/\y in {2/1,3/1,3/2,4/1,5.5/1,5.5/2,6.5/1,6.5/2,9/1}
  \fill[quotientorange] (\x,0.3*\y+0.12) circle (3pt);
\end{scope}
\end{tikzpicture}
\caption{Every possible mean of two draws from $\{2,4,9\}$. The means are centred
on the same value as the population but are less spread out: variance $4.33$
instead of $8.67$.}
\label{fig:sd-tiny}
\end{figure}

\section{The sampling distribution of the mean, by simulation}
\label{sec:sd-sim}

Real populations are too large to list. We simulate instead. Our population is
request latency, modelled as a fixed 20 ms floor plus an exponential delay with
mean 30 ms. This population is strongly right-skewed: a long tail of slow
requests stretches to the right. Its \textbf{skewness}, a unitless measure of
that lopsidedness (0 for a symmetric shape, positive for a long right tail;
Chapter 3), is 2. It has
\[
\mu=50\text{ ms},\qquad \sigma=30\text{ ms},\qquad \text{median}=20+30\ln 2\approx 40.8\text{ ms}.
\]
For each $n\in\{1,4,16,64\}$ we drew $10{,}000$ samples of size $n$, computed each
sample's mean, and histogrammed the $10{,}000$ means (\texttt{numpy} seed 6).

\begin{lstlisting}[language=Python]
import numpy as np
rng = np.random.default_rng(6)
for n in [1, 4, 16, 64]:
    samples = 20 + rng.exponential(30, size=(10_000, n))
    means = samples.mean(axis=1)          # one mean per simulated study
    print(n, means.mean().round(2), means.std(ddof=1).round(2))
\end{lstlisting}

\begin{figure}[H]
\centering
\begin{tikzpicture}
\pgfplotsset{sdpanel/.style={width=6.6cm, height=4.4cm, ymin=0, axis lines=left,
  tick label style={font=\scriptsize}, title style={font=\small},
  xlabel style={font=\scriptsize}, xlabel={mean latency (ms)}}}
\begin{axis}[sdpanel, at={(0,0)}, xmin=20, xmax=200, ymax=0.034,
  title={$n=1$ (the population)}, title style={font=\small}, scaled y ticks=false,
  yticklabel style={/pgf/number format/fixed, /pgf/number format/precision=2}]
\addplot[ybar interval, fill=histfill, draw=white] coordinates {(20,0.02843) (30,0.02067) (40,0.01476) (50,0.01040) (60,0.00697) (70,0.00581) (80,0.00366) (90,0.00271) (100,0.00182) (110,0.00137) (120,0.00110) (130,0.00071) (140,0.00053) (150,0.00028) (160,0.00025) (170,0.00010) (180,0.00010) (190,0.00011) (200,0)};
\addplot[densityred, thick, domain=20:200, samples=80] {exp(-((x-50)^2)/(2*900))/(30*2.5066)};
\draw[black, dashed] (axis cs:50,0) -- (axis cs:50,0.034);
\end{axis}
\begin{axis}[sdpanel, at={(6.3cm,0)}, xmin=20, xmax=120, ymax=0.034,
  title={$n=4$}, title style={font=\small}, scaled y ticks=false,
  yticklabel style={/pgf/number format/fixed, /pgf/number format/precision=2}]
\addplot[ybar interval, fill=histfill, draw=white] coordinates {(20,0.00082) (25,0.00852) (30,0.01884) (35,0.02718) (40,0.02882) (45,0.02766) (50,0.02462) (55,0.01912) (60,0.01420) (65,0.01082) (70,0.00676) (75,0.00418) (80,0.00300) (85,0.00166) (90,0.00130) (95,0.00074) (100,0.00062) (105,0.00052) (110,0.00034) (115,0.00012) (120,0)};
\addplot[densityred, thick, domain=20:120, samples=80] {exp(-((x-50)^2)/(2*225))/(15*2.5066)};
\draw[black, dashed] (axis cs:50,0) -- (axis cs:50,0.034);
\end{axis}
\begin{axis}[sdpanel, at={(0,-5.1cm)}, xmin=25, xmax=85, ymax=0.06,
  title={$n=16$}, title style={font=\small}, scaled y ticks=false,
  yticklabel style={/pgf/number format/fixed, /pgf/number format/precision=2}]
\addplot[ybar interval, fill=histfill, draw=white] coordinates {(25,0.00000) (28,0.00010) (31,0.00157) (34,0.00703) (37,0.01530) (40,0.03363) (43,0.04880) (46,0.05337) (49,0.05207) (52,0.04167) (55,0.03217) (58,0.02087) (61,0.01303) (64,0.00707) (67,0.00357) (70,0.00173) (73,0.00053) (76,0.00057) (79,0.00013) (82,0.00013) (85,0)};
\addplot[densityred, thick, domain=25:85, samples=80] {exp(-((x-50)^2)/(2*56.25))/(7.5*2.5066)};
\draw[black, dashed] (axis cs:50,0) -- (axis cs:50,0.06);
\end{axis}
\begin{axis}[sdpanel, at={(6.3cm,-5.1cm)}, xmin=35, xmax=65, ymax=0.115,
  title={$n=64$}, title style={font=\small}, scaled y ticks=false,
  yticklabel style={/pgf/number format/fixed, /pgf/number format/precision=2}]
\addplot[ybar interval, fill=histfill, draw=white] coordinates {(35,0.00000) (36.5,0.00013) (38,0.00013) (39.5,0.00213) (41,0.00773) (42.5,0.01867) (44,0.04607) (45.5,0.07187) (47,0.09820) (48.5,0.10347) (50,0.09860) (51.5,0.08180) (53,0.05773) (54.5,0.03880) (56,0.02193) (57.5,0.01100) (59,0.00480) (60.5,0.00213) (62,0.00080) (63.5,0.00047) (65,0)};
\addplot[densityred, thick, domain=35:65, samples=80] {exp(-((x-50)^2)/(2*14.0625))/(3.75*2.5066)};
\draw[black, dashed] (axis cs:50,0) -- (axis cs:50,0.115);
\end{axis}
\end{tikzpicture}
\caption{Histograms of $10{,}000$ simulated sample means from the skewed latency
population, for four sample sizes. Note the shrinking $x$-ranges. The red curve
is $\Normal(50, 30^2/n)$. As $n$ grows the histogram narrows around $\mu=50$
(dashed) and its skew fades until the Normal curve fits well.}
\label{fig:sd-grid}
\end{figure}

The simulation output:

\begin{center}
\begin{tabular}{rcccc}
\toprule
$n$ & mean of the $\xbar$'s & SD of the $\xbar$'s & $30/\sqrt n$ & skewness of the $\xbar$'s\\
\midrule
1  & 49.56 & 29.38 & 30.00 & 1.98\\
4  & 50.12 & 15.08 & 15.00 & 1.04\\
16 & 50.02 & 7.47  & 7.50  & 0.48\\
64 & 49.98 & 3.76  & 3.75  & 0.29\\
\bottomrule
\end{tabular}
\end{center}

Three patterns stand out. The centre stays at $\mu=50$. The spread follows
$30/\sqrt n$ almost exactly: quadrupling $n$ halves it. The shape becomes more
symmetric; the skewness falls roughly like $2/\sqrt n$ (theory gives
$2,1,0.5,0.25$). The next three sections explain each pattern.

\section{Deriving the centre and the spread}
\label{sec:sd-derive}

Let $X_1,\dots,X_n\iid$ some population with mean $\mu$ and variance $\sigma^2$.
``iid'' means independent and identically distributed: each draw comes from the
same population and does not influence the others. The sample mean is
$\Xbar=\frac1n(X_1+\dots+X_n)$.

\begin{theorem}[Mean and variance of the sample mean]
\label{thm:sd-xbar}
\[
\E[\Xbar]=\mu,\qquad \Var(\Xbar)=\frac{\sigma^2}{n},\qquad
\SE(\Xbar)=\SD(\Xbar)=\frac{\sigma}{\sqrt n}.
\]
\end{theorem}

\begin{whybox}
\textbf{Centre.} By linearity (rule 1),
\[
\E[\Xbar]=\frac1n\bigl(\E[X_1]+\dots+\E[X_n]\bigr)=\frac1n\cdot n\mu=\mu.
\]
No independence was needed. $\Xbar$ is \emph{unbiased}: on average it hits the
target.

\textbf{Spread.} Variances of independent variables add (rule 3), so the
\emph{sum} has variance
\[
\Var(X_1+\dots+X_n)=n\sigma^2 .
\]
Dividing by $n$ multiplies variance by $1/n^2$ (rule 2):
\[
\Var(\Xbar)=\frac{1}{n^2}\cdot n\sigma^2=\frac{\sigma^2}{n}.
\]
The sum's variance grows like $n$, but the division shrinks it by $n^2$. The net
effect is $1/n$ on the variance, hence $1/\sqrt n$ on the SD.
\end{whybox}

Check with the tiny population: $\sigma^2/n=(26/3)/2=13/3\approx 4.333$. Check with
the simulation: $30/\sqrt{16}=7.5$, observed $7.47$. (The pause question: with
$n=3$ the variance would be $(26/3)/3=26/9\approx 2.889$.)

\subsection{Why $\sqrt n$, and what it costs you}

The standard error falls like $1/\sqrt n$, not $1/n$. To halve the error you need
four times the data. To cut it by 10 you need 100 times the data.

\begin{figure}[H]
\centering
\begin{tikzpicture}
\begin{axis}[width=11cm, height=5.6cm, xmin=0, xmax=100, ymin=0, ymax=32,
  xlabel={sample size $n$}, ylabel={$\SE=30/\sqrt n$ (ms)}, axis lines=left,
  xtick={4,16,36,64,100}, ytick={0,3.75,7.5,15,30},
  yticklabel style={/pgf/number format/fixed}]
\addplot[distblue, very thick, domain=1:100, samples=120] {30/sqrt(x)};
\foreach \n/\s in {4/15, 16/7.5, 64/3.75} {
  \edef\temp{\noexpand\draw[gray, dashed] (axis cs:\n,0) -- (axis cs:\n,\s) -- (axis cs:0,\s);
    \noexpand\fill[quotientorange] (axis cs:\n,\s) circle (2.5pt);}
  \temp
}
\fill[quotientorange] (axis cs:1,30) circle (2.5pt);
\node[anchor=west, font=\small] at (axis cs:6,16) {$n=4$: 15 ms};
\node[anchor=west, font=\small] at (axis cs:18,9.5) {$\times 4$ data: 7.5 ms};
\node[anchor=west, font=\small] at (axis cs:66,6) {$\times 4$ again: 3.75 ms};
\end{axis}
\end{tikzpicture}
\caption{Standard error of the mean against $n$ for $\sigma=30$ ms. Each
quadrupling of the sample size halves the error. The curve flattens fast:
the first few observations buy a lot of precision, later ones buy little.}
\label{fig:sd-sqrtn}
\end{figure}

\begin{intuitionbox}
Why not $1/n$? Errors in independent observations point in random directions.
Adding $n$ of them is like a random walk: the total drifts about $\sqrt n$ steps
from zero, not $n$ steps, because many errors cancel. Dividing that $\sqrt n$
drift by $n$ leaves $1/\sqrt n$.
\end{intuitionbox}

\begin{recoverybox}
Pause and think: an A/B dashboard shows a conversion-rate SE of 0.4 points after
one week. Roughly what will it show after four weeks of the same traffic?
(Half: 0.2 points, since four times the data halves the SE.)
\end{recoverybox}

\begin{example}[Planning a sample size]
A benchmark of 25 runs gives SE $=\sigma/5$. The team wants the SE cut to a third.
They need $n=25\times 3^2=225$ runs. Generally, $n_{\text{new}}=n_{\text{old}}\cdot
(\SE_{\text{old}}/\SE_{\text{new}})^2$.
\end{example}

\section{Standard deviation versus standard error}
\label{sec:sd-sdvsse}

These two numbers sound alike and answer different questions.

\begin{itemize}
  \item \textbf{SD} ($\sigma$, estimated by $s$): how much \emph{individual
  observations} vary. A property of the population. Does not shrink with $n$.
  \item \textbf{SE} ($\sigma/\sqrt n$): how much the \emph{statistic} would vary
  across repeated studies. A property of your estimation procedure. Shrinks with
  $n$.
\end{itemize}

\begin{figure}[H]
\centering
\begin{tikzpicture}
\begin{axis}[width=12cm, height=5.6cm, xmin=0, xmax=160, ymin=0, ymax=0.068,
  xlabel={latency (ms)}, ylabel={density}, axis lines=left,
  scaled y ticks=false, yticklabel style={/pgf/number format/fixed, /pgf/number format/precision=2},
  legend style={font=\small, at={(0.98,0.95)}, anchor=north east}]
\addplot[distblue, very thick, domain=20:160, samples=100] {exp(-(x-20)/30)/30};
\addlegendentry{one request: SD $=30$ ms}
\addplot[densityred, very thick, domain=20:80, samples=100] {exp(-((x-50)^2)/(2*56.25))/(7.5*2.5066)};
\addlegendentry{mean of 16 requests: SE $=7.5$ ms}
\draw[distblue, thick] (axis cs:20,0) -- (axis cs:20,0.0333);
\draw[gray, dashed] (axis cs:50,0) -- (axis cs:50,0.056);
\node[font=\small, anchor=south] at (axis cs:50,0.056) {$\mu=50$};
\end{axis}
\end{tikzpicture}
\caption{The population of single latencies (blue) and the sampling distribution
of the mean of 16 of them (red, Normal approximation). Collecting more data
does not change the blue curve. It only squeezes the red one.}
\label{fig:sd-sdse}
\end{figure}

\begin{trapbox}
\textbf{Error bars of $\pm$SE presented as if they described individuals.} A paper
reports ``mean recovery time $12.0 \pm 0.4$ days (SE), $n=400$.'' The SD is
$0.4\sqrt{400}=8$ days. Individual patients vary by about 8 days, not 0.4. Always
check which one a $\pm$ means. Rule of thumb: if you want to predict \emph{one
new observation}, use SD; if you want to know \emph{how well you pinned down the
mean}, use SE.

The reverse mistake also happens: reporting SD bars on a comparison of two means
makes a solid difference look like noise.
\end{trapbox}

\section{The sampling distribution of a proportion}
\label{sec:sd-prop}

A proportion is a mean in disguise. Code each visitor as $X_i=1$ if they convert
and $0$ otherwise. Then the sample proportion is
\[
\hat p=\frac{\text{number of 1s}}{n}=\frac{X_1+\dots+X_n}{n}=\Xbar .
\]
Each $X_i$ is Bernoulli with $\E[X_i]=p$ and
\[
\Var(X_i)=\E[X_i^2]-p^2=p-p^2=p(1-p),
\]
using $X_i^2=X_i$ for a 0/1 variable. Theorem~\ref{thm:sd-xbar} then gives
everything for free.

\begin{theorem}[Sample proportion]
For $n$ independent 0/1 observations with success probability $p$,
\[
\E[\hat p]=p,\qquad \SE(\hat p)=\sqrt{\frac{p(1-p)}{n}}.
\]
\end{theorem}

\begin{example}[Conversion rate]
A landing page converts at $p=0.10$. With $n=400$ visitors,
$\SE(\hat p)=\sqrt{0.1\cdot 0.9/400}=\sqrt{0.000225}=0.015$, i.e.\ 1.5 percentage
points. Observed rates of 8.5\% or 11.5\% are one SE from the truth and entirely
routine.
\end{example}

The numerator $p(1-p)$ is largest at $p=0.5$, where it equals $0.25$
(Figure~\ref{fig:sd-pq}). So $\SE(\hat p)\le 0.5/\sqrt n$ for every $p$. A poll of
$n=1000$ has $\SE\le 0.5/\sqrt{1000}=0.0158$, and twice that, $0.032$, is the
``$\pm 3$ points'' margin printed under news polls (Chapter 7 explains the factor
of about 2).

\begin{figure}[H]
\centering
\begin{tikzpicture}
\begin{axis}[width=10cm, height=5.2cm, xmin=0, xmax=1, ymin=0, ymax=0.58,
  xlabel={true proportion $p$}, ylabel={$\sqrt{p(1-p)}$}, axis lines=left,
  xtick={0,0.1,0.25,0.5,0.75,0.9,1}, ytick={0,0.3,0.433,0.5}]
\addplot[structurepurple, very thick, domain=0:1, samples=200] {sqrt(x*(1-x))};
\draw[gray, dashed] (axis cs:0.5,0) -- (axis cs:0.5,0.5) -- (axis cs:0,0.5);
\draw[gray, dashed] (axis cs:0.1,0) -- (axis cs:0.1,0.3) -- (axis cs:0,0.3);
\fill[quotientorange] (axis cs:0.5,0.5) circle (2.5pt);
\fill[quotientorange] (axis cs:0.1,0.3) circle (2.5pt);
\node[anchor=south, font=\small] at (axis cs:0.5,0.51) {worst case $0.5$};
\node[anchor=west, font=\small] at (axis cs:0.12,0.25) {$p=0.1$: $0.3$};
\end{axis}
\end{tikzpicture}
\caption{The per-observation SD of a 0/1 variable, $\sqrt{p(1-p)}$. It peaks at
$p=0.5$ and is symmetric. Divide by $\sqrt n$ to get $\SE(\hat p)$. Using $0.5$
gives a conservative SE when $p$ is unknown.}
\label{fig:sd-pq}
\end{figure}

When $p$ is near 0 or 1 the sampling distribution of $\hat p$ is lopsided for small
$n$, because $\hat p$ cannot go below 0. Figure~\ref{fig:sd-phat} shows $p=0.05$.
The usual rule is that the Normal approximation is fine when both $np\ge 10$ and
$n(1-p)\ge 10$.

\begin{figure}[H]
\centering
\begin{tikzpicture}
\begin{axis}[at={(0,0)}, width=6.6cm, height=4.6cm, ybar, bar width=5pt,
  xmin=-0.02, xmax=0.3, ymin=0, ymax=0.45, axis lines=left,
  title={$n=20$, $np=1$}, title style={font=\small}, xlabel={\small $\hat p$}, ylabel={\small probability},
  tick label style={font=\scriptsize}, xtick={0,0.05,0.1,0.2,0.3},
  xticklabel style={/pgf/number format/fixed}]
\addplot[fill=histfill, draw=distblue] coordinates {(0.000,0.3585) (0.050,0.3774) (0.100,0.1887) (0.150,0.0596) (0.200,0.0133) (0.250,0.0022)};
\addplot[densityred, thick, sharp plot, domain=-0.02:0.3, samples=100] {0.05*exp(-((x-0.05)^2)/(2*0.002375))/(0.04873*2.5066)};
\end{axis}
\begin{axis}[at={(6.4cm,0)}, width=6.6cm, height=4.6cm, ybar, bar width=1.2pt,
  xmin=-0.02, xmax=0.3, ymin=0, ymax=0.15, axis lines=left,
  title={$n=200$, $np=10$}, title style={font=\small}, xlabel={\small $\hat p$},
  tick label style={font=\scriptsize}, xtick={0,0.05,0.1,0.2,0.3},
  xticklabel style={/pgf/number format/fixed}, scaled y ticks=false,
  yticklabel style={/pgf/number format/fixed}]
\addplot[fill=histfill, draw=distblue] coordinates {(0.010,0.0019) (0.015,0.0067) (0.020,0.0174) (0.025,0.0359) (0.030,0.0614) (0.035,0.0896) (0.040,0.1137) (0.045,0.1277) (0.050,0.1284) (0.055,0.1167) (0.060,0.0967) (0.065,0.0736) (0.070,0.0518) (0.075,0.0338) (0.080,0.0206) (0.085,0.0117) (0.090,0.0063) (0.095,0.0032) (0.100,0.0015) (0.105,0.0007)};
\addplot[densityred, thick, sharp plot, domain=-0.02:0.3, samples=150] {0.005*exp(-((x-0.05)^2)/(2*0.0002375))/(0.01541*2.5066)};
\end{axis}
\end{tikzpicture}
\caption{Exact (binomial) sampling distributions of $\hat p$ when $p=0.05$, with
the Normal approximation in red. At $n=20$ the Normal curve puts real
probability on impossible negative proportions. At $n=200$ it fits.}
\label{fig:sd-phat}
\end{figure}

\section{The shape: Central Limit Theorem and Law of Large Numbers}
\label{sec:sd-clt}

Theorem~\ref{thm:sd-xbar} gives the centre and spread of $\Xbar$ for any
population. The shape comes from the Central Limit Theorem (CLT).

\begin{theorem}[Central Limit Theorem]
\label{thm:sd-clt}
If $X_1,\dots,X_n\iid$ a population with mean $\mu$ and finite variance
$\sigma^2$, then for large $n$
\[
\Xbar\ \overset{\text{approx}}{\sim}\ \Normal\!\left(\mu,\ \frac{\sigma^2}{n}\right),
\qquad\text{equivalently}\qquad
Z=\frac{\Xbar-\mu}{\sigma/\sqrt n}\ \overset{\text{approx}}{\sim}\ \Normal(0,1).
\]
\end{theorem}

The proof is in \emph{Probability From the Ground Up}. What matters here is what it
buys: once you know $\mu$ and $\SE$, you know the whole sampling distribution, so
you know, for example, that $\Xbar$ lands within $2\SE$ of $\mu$ about 95\% of the
time. Chapter 7 turns this sentence into confidence intervals.

\begin{historybox}
Abraham de Moivre, a French mathematician working in London, was stuck in 1733 on
a gambler's computation: the probability of, say, between 3400 and 3600 heads in
6000 coin tosses needs a sum of thousands of enormous binomial coefficients. He
found that the binomial probabilities trace out a bell-shaped curve, the curve we
now call the Normal density, so the sum could be replaced by an area under it.
Pierre-Simon Laplace extended the result around 1810 from coin tosses to sums of
almost any independent quantities, and used it to put error bars on astronomical
averages. That generalisation is the Central Limit Theorem. It is the reason one
bell curve and one number, the standard error, are enough to describe the
uncertainty of almost any average.
\end{historybox}

\subsection{When is $n$ ``large enough''?}

There is no universal number. ``$n\ge 30$'' is folklore that works for mildly
skewed data and fails for heavily skewed data. The governing quantity is the
population's skewness $\gamma$. The skewness of $\Xbar$ is exactly
\[
\gamma_{\Xbar}=\frac{\gamma}{\sqrt n}.
\]
(Our latency population has $\gamma=2$; the simulated skewness of the means was
$1.98, 1.04, 0.48, 0.29$ against theory $2, 1, 0.5, 0.25$.) A skewness below about
$0.3$ is usually close enough to Normal for intervals and tests, which needs
$n\ge(\gamma/0.3)^2$: at least 45 for $\gamma=2$, but about 1100 for $\gamma=10$.

\begin{figure}[H]
\centering
\begin{tikzpicture}
\pgfplotsset{cltpanel/.style={width=4.8cm, height=4.2cm, ymin=0, axis lines=left,
  tick label style={font=\scriptsize}, title style={font=\small}}}
\begin{axis}[cltpanel, at={(0,0)}, xmin=-1, xmax=3, ymax=0.85, title={$n=2$}, title style={font=\small}]
\addplot[distblue, thick, smooth] coordinates {(0,0) (0.05,0.1810) (0.1,0.3275) (0.2,0.5363) (0.3,0.6586) (0.4,0.7189) (0.5,0.7358) (0.6,0.7229) (0.7,0.6905) (0.8,0.6461) (0.9,0.5951) (1.0,0.5413) (1.1,0.4875) (1.2,0.4354) (1.3,0.3862) (1.4,0.3405) (1.5,0.2987) (1.6,0.2609) (1.7,0.2269) (1.8,0.1967) (1.9,0.1700) (2.0,0.1465) (2.2,0.1080) (2.4,0.0790) (2.6,0.0574) (2.8,0.0414) (3.0,0.0297)};
\addplot[densityred, thick, dashed, domain=-1:3, samples=80] {exp(-((x-1)^2)/(2*0.5))/(0.7071*2.5066)};
\end{axis}
\begin{axis}[cltpanel, at={(4.25cm,0)}, xmin=0, xmax=2.4, ymax=1.5, title={$n=10$}, title style={font=\small}]
\addplot[distblue, thick, smooth] coordinates {(0.2,0.0019) (0.3,0.0270) (0.4,0.1323) (0.5,0.3627) (0.6,0.6884) (0.7,1.0140) (0.8,1.2408) (0.9,1.3176) (1.0,1.2511) (1.1,1.0853) (1.2,0.8736) (1.3,0.6605) (1.4,0.4734) (1.5,0.3241) (1.6,0.2131) (1.7,0.1353) (1.8,0.0833) (1.9,0.0498) (2.0,0.0291) (2.2,0.0093) (2.4,0.0027)};
\addplot[densityred, thick, dashed, domain=0:2.4, samples=80] {exp(-((x-1)^2)/(2*0.1))/(0.3162*2.5066)};
\end{axis}
\begin{axis}[cltpanel, at={(8.5cm,0)}, xmin=0.5, xmax=1.6, ymax=2.9, title={$n=40$}, title style={font=\small}]
\addplot[distblue, thick, smooth] coordinates {(0.5,0.0022) (0.55,0.0124) (0.6,0.0498) (0.65,0.1530) (0.7,0.3727) (0.75,0.7437) (0.8,1.2471) (0.85,1.7953) (0.9,2.2577) (0.95,2.5168) (0.975,2.5498) (1.0,2.5179) (1.05,2.2847) (1.1,1.8975) (1.15,1.4537) (1.2,1.0345) (1.25,0.6880) (1.3,0.4298) (1.35,0.2535) (1.4,0.1417) (1.45,0.0754) (1.5,0.0383) (1.55,0.0186) (1.6,0.0087)};
\addplot[densityred, thick, dashed, domain=0.5:1.6, samples=80] {exp(-((x-1)^2)/(2*0.025))/(0.1581*2.5066)};
\end{axis}
\end{tikzpicture}
\caption{Exact density of the mean of $n$ exponential observations with mean 1
(blue; skewness $2/\sqrt n$) against the CLT's Normal approximation (red, dashed).
At $n=2$ the Normal curve even puts mass below zero. At $n=10$ the peak is
visibly off-centre. At $n=40$ the two nearly coincide.}
\label{fig:sd-clt}
\end{figure}

\begin{trapbox}
The CLT needs a finite variance. For extremely heavy-tailed data (some file
sizes, some incomes; mathematically, Pareto or ``power-law'' tails, where the
chance of exceeding $x$ falls only like $1/x^{\alpha}$ with $\alpha<2$) the
sample mean converges
painfully slowly or not at all, and a few huge values dominate every sample. For
such data report medians or quantiles (Chapter 3), or use the bootstrap
(Chapter 8) with care.
\end{trapbox}

\subsection{The Law of Large Numbers}

The CLT describes the \emph{fluctuations} of $\Xbar$. The Law of Large Numbers
(LLN) is the cruder statement that those fluctuations vanish:
\[
\Xbar_n\to\mu\quad\text{as } n\to\infty .
\]
It follows directly from $\Var(\Xbar)=\sigma^2/n\to 0$: a quantity whose spread
around $\mu$ shrinks to zero must end up at $\mu$. (The formal step is Chebyshev's
inequality, which bounds the chance of landing $k$ SDs away by $1/k^2$; see
\emph{Probability From the Ground Up}.) Figure~\ref{fig:sd-lln} shows a single
running average of latency draws settling down inside the funnel
$\mu\pm 2\sigma/\sqrt n$.

\begin{figure}[H]
\centering
\begin{tikzpicture}
\begin{axis}[width=12cm, height=5.6cm, xmin=0, xmax=500, ymin=20, ymax=90,
  xlabel={number of requests $n$}, ylabel={running mean $\xbar_n$ (ms)}, axis lines=left,
  legend style={font=\small, at={(0.98,0.97)}, anchor=north east}]
\addplot[name path=up, gray, thin, domain=4:500, samples=150, forget plot] {50+60/sqrt(x)};
\addplot[name path=lo, gray, thin, domain=4:500, samples=150, forget plot] {50-60/sqrt(x)};
\addplot[gray!20, area legend] fill between[of=up and lo];
\addlegendentry{$\mu\pm 2\sigma/\sqrt n$}
\addplot[distblue, thick] coordinates {(1,59.47) (2,69.93) (3,65.73) (4,60.63) (5,56.03) (6,50.48) (7,55.16) (8,53.33) (9,51.26) (10,57.25) (11,55.82) (12,54.95) (13,57.45) (14,58.21) (15,56.41) (16,56.21) (17,54.53) (18,54.77) (19,53.96) (20,52.52) (30,47.80) (40,48.84) (50,50.27) (60,48.26) (70,49.04) (80,49.16) (90,50.58) (100,50.83) (110,49.59) (120,49.32) (130,48.99) (140,48.74) (150,48.46) (160,48.44) (170,48.57) (180,49.01) (190,48.22) (200,48.00) (210,47.93) (220,47.45) (230,48.04) (240,47.65) (250,47.75) (260,47.74) (270,48.31) (280,48.26) (290,48.43) (300,47.96) (310,48.41) (320,48.30) (330,48.35) (340,48.81) (350,48.56) (360,48.67) (370,48.88) (380,48.69) (390,48.87) (400,48.71) (410,48.69) (420,48.46) (430,48.27) (440,48.05) (450,48.51) (460,48.61) (470,48.51) (480,48.43) (490,48.48) (500,48.72)};
\addlegendentry{running mean}
\addplot[densityred, dashed, domain=0:500] {50};
\addlegendentry{$\mu=50$}
\end{axis}
\end{tikzpicture}
\caption{One stream of 500 simulated latencies. The running mean wanders early
(it is 69.9 ms after two requests) and then settles near $\mu=50$, ending at
48.7 ms. The shaded funnel narrows like $1/\sqrt n$: the LLN says the mean
converges, and the SE says how fast.}
\label{fig:sd-lln}
\end{figure}

\begin{trapbox}
The LLN does not say that luck ``evens out'' by compensation. After a run of slow
requests, the next request is no more likely to be fast. The early excess is not
cancelled; it is \emph{diluted} by the growing denominator.
\end{trapbox}

\section{Estimated standard error}
\label{sec:sd-estse}

The formula $\sigma/\sqrt n$ contains $\sigma$, the population SD, which you do not
know. Replace it with the sample SD $s$ (Chapter 2, with the $n-1$ divisor):
\[
\widehat{\SE}(\Xbar)=\frac{s}{\sqrt n},\qquad
\widehat{\SE}(\hat p)=\sqrt{\frac{\hat p(1-\hat p)}{n}} .
\]
In practice ``the SE'' almost always means this estimated version.

\begin{example}[Eight latencies]
Measured latencies (ms): $42, 35, 61, 48, 120, 39, 55, 44$. The sum is 444, so
$\xbar=55.5$. The squared deviations sum to $5254$, so
$s=\sqrt{5254/7}=\sqrt{750.57}\approx 27.40$ ms and
\[
\widehat{\SE}=\frac{27.40}{\sqrt 8}=\frac{27.40}{2.828}\approx 9.69\text{ ms}.
\]
Report: mean latency $55.5$ ms, SE $9.7$ ms. With only 8 skewed observations, the
Normal shape for $\Xbar$ is doubtful and the plug-in $s$ is itself noisy;
Chapter 7 handles the second issue with the $t$ distribution.
\end{example}

\begin{example}[Estimated SE of a proportion]
48 of 400 visitors converted: $\hat p=0.12$ and
$\widehat{\SE}=\sqrt{0.12\cdot 0.88/400}=\sqrt{0.000264}\approx 0.0162$.
\end{example}

\begin{lstlisting}[language=Python]
import numpy as np
from scipy import stats
x = np.array([42, 35, 61, 48, 120, 39, 55, 44])
print(x.std(ddof=1) / np.sqrt(len(x)))   # 9.686
print(stats.sem(x))                       # same: scipy uses ddof=1
\end{lstlisting}

\section{Standard error of a difference}
\label{sec:sd-diff}

A/B tests, before/after comparisons and treatment/control studies all estimate a
\emph{difference} $\Xbar_A-\Xbar_B$. If the two groups are independent, rule 3
(variances add, even for a difference) gives
\[
\Var(\Xbar_A-\Xbar_B)=\Var(\Xbar_A)+\Var(\Xbar_B),
\qquad
\SE(\Xbar_A-\Xbar_B)=\sqrt{\SE_A^2+\SE_B^2}=\sqrt{\frac{\sigma_A^2}{n_A}+\frac{\sigma_B^2}{n_B}} .
\]

\begin{example}[A/B latency test]
Variant A: $n_A=100$, $s_A=30$ ms, so $\SE_A=3.00$. Variant B: $n_B=150$, $s_B=40$
ms, so $\SE_B=40/\sqrt{150}\approx 3.27$.
\[
\SE(\xbar_A-\xbar_B)=\sqrt{\frac{900}{100}+\frac{1600}{150}}=\sqrt{9+10.67}=\sqrt{19.67}\approx 4.43\text{ ms}.
\]
Adding the SEs would give $6.27$ ms, an overstatement of 41\%.
\end{example}

\begin{figure}[H]
\centering
\begin{tikzpicture}[scale=1.0]
\coordinate (O) at (0,0); \coordinate (A) at (3.0,0); \coordinate (B) at (3.0,3.27);
\draw[very thick, distblue] (O) -- (A) node[midway, below, font=\small] {$\SE_A=3.00$};
\draw[very thick, targetgreen] (A) -- (B) node[midway, right, font=\small] {$\SE_B=3.27$};
\draw[very thick, densityred] (O) -- (B) node[midway, above left, font=\small] {$\SE_{A-B}=4.43$};
\draw (2.75,0) -- (2.75,0.25) -- (3.0,0.25);
\begin{scope}[xshift=6cm]
\draw[very thick, distblue] (0,0) -- (3.0,0);
\draw[very thick, targetgreen] (3.0,0) -- (6.27,0);
\draw[trapred, thick, decorate, decoration={brace, amplitude=6pt, mirror}] (0,-0.15) -- (6.27,-0.15)
  node[midway, below=7pt, font=\small, trapred] {wrong: $3.00+3.27=6.27$};
\end{scope}
\end{tikzpicture}
\caption{Standard errors of independent estimates combine like the sides of a
right triangle (Pythagoras), not like segments laid end to end. The hypotenuse
is always shorter than the sum of the legs.}
\label{fig:sd-pythag}
\end{figure}

\begin{whybox}
Why does the difference not have a \emph{smaller} variance? Subtracting $\Xbar_B$
means multiplying it by $a=-1$, and rule 2 says that multiplies its variance by
$a^2=+1$. Noise in B pushes the difference around just as much whichever sign it
carries. Independence then makes the two variances add.
\end{whybox}

\begin{trapbox}
The formula needs independent groups. In a paired design (the same 50 servers
measured before and after a patch) compute the 50 differences first and use
$s_{\text{diff}}/\sqrt{50}$. Treating paired data as independent usually
\emph{overstates} the SE because it ignores the positive correlation between the
two measurements on the same server (Chapter 9).
\end{trapbox}

\section{Statistics beyond the mean: the median}
\label{sec:sd-median}

Every statistic has a sampling distribution and an SE, not only the mean. The
median is the natural next case. We simulated $10{,}000$ samples of $n=25$ heights
from $\Normal(170, 10^2)$ and recorded both the mean and the median of each.

\begin{figure}[H]
\centering
\begin{tikzpicture}
\begin{axis}[width=11.5cm, height=5.6cm, xmin=162, xmax=178, ymin=0, ymax=0.28,
  xlabel={statistic (cm)}, ylabel={density}, axis lines=left,
  scaled y ticks=false, yticklabel style={/pgf/number format/fixed, /pgf/number format/precision=2},
  legend style={font=\small, at={(0.98,0.97)}, anchor=north east}]
\addplot[ybar interval, fill=histfill, draw=white] coordinates {(162,0.00025) (162.8,0.00087) (163.6,0.00287) (164.4,0.00662) (165.2,0.02050) (166,0.03925) (166.8,0.07862) (167.6,0.12037) (168.4,0.16625) (169.2,0.19837) (170,0.19087) (170.8,0.16200) (171.6,0.12050) (172.4,0.07112) (173.2,0.04250) (174,0.01762) (174.8,0.00750) (175.6,0.00300) (176.4,0.00087) (177.2,0.00000) (178,0)};
\addlegendentry{sample means (SD $2.02$)}
\addplot[const plot mark left, bootorange, very thick] coordinates {(162,0.00137) (162.8,0.00512) (163.6,0.00900) (164.4,0.01675) (165.2,0.03675) (166,0.05837) (166.8,0.08537) (167.6,0.11512) (168.4,0.14512) (169.2,0.16425) (170,0.15500) (170.8,0.13637) (171.6,0.11212) (172.4,0.08400) (173.2,0.05725) (174,0.03275) (174.8,0.01950) (175.6,0.00800) (176.4,0.00400) (177.2,0.00162) (178,0.00162)};
\addlegendentry{sample medians (SD $2.50$)}
\end{axis}
\end{tikzpicture}
\caption{Sampling distributions of the mean and median for $n=25$ draws from
$\Normal(170,10^2)$. Both are centred on 170. The median's histogram is lower and
wider: for Normal data the median is a noisier estimate of the centre than the
mean.}
\label{fig:sd-medmean}
\end{figure}

The mean's SD matches $10/\sqrt{25}=2$. The median's SD, $2.50$, matches the
large-sample formula below: $1.2533\times 2=2.51$.

\begin{theorem}[Large-sample SE of the median]
\label{thm:sd-median}
If the population has density $f$ and median $m$ with $f(m)>0$, then for large
$n$ the sample median is approximately Normal with mean $m$ and
\[
\SE(\text{median})\approx\frac{1}{2f(m)\sqrt n}.
\]
For a Normal population $f(m)=1/(\sigma\sqrt{2\pi})$, which gives
\[
\SE(\text{median})\approx\frac{\sigma\sqrt{2\pi}}{2\sqrt n}=\sqrt{\frac{\pi}{2}}\,\frac{\sigma}{\sqrt n}\approx 1.2533\,\frac{\sigma}{\sqrt n}.
\]
\end{theorem}

\begin{whybox}
The median is the point with half the sample below it. Count how many
observations fall below the true median $m$. That count is binomial with $p=1/2$,
so the \emph{fraction} below $m$ has SE $\sqrt{0.25/n}=0.5/\sqrt n$. This wobble
lives on the vertical (probability) axis of the CDF. To turn it into a horizontal
wobble in the median, divide by the slope of the CDF at $m$, which is the density
$f(m)$. Result: $0.5/(f(m)\sqrt n)=1/(2f(m)\sqrt n)$. A steep CDF (lots of data
near the median) pins the median down; a flat one lets it slide
(Figure~\ref{fig:sd-cdfmed}).
\end{whybox}

\begin{figure}[H]
\centering
\begin{tikzpicture}
\begin{axis}[width=11cm, height=6cm, xmin=-4, xmax=4, ymin=0, ymax=1,
  xlabel={value}, ylabel={CDF $F$}, axis lines=left, xtick={0}, xticklabels={$m$},
  ytick={0.3,0.5,0.7}, yticklabels={$\frac12-\delta$, $\frac12$, $\frac12+\delta$},
  legend style={font=\small, at={(0.03,0.97)}, anchor=north west}]
\fill[quotientorange!20] (axis cs:-4,0.3) rectangle (axis cs:4,0.7);
\addplot[distblue, very thick, domain=-4:4, samples=100] {1/(1+exp(-2*x))};
\addlegendentry{steep CDF (large $f(m)$)}
\addplot[densityred, very thick, domain=-4:4, samples=100] {1/(1+exp(-0.6*x))};
\addlegendentry{flat CDF (small $f(m)$)}
\fill[distblue, opacity=0.35] (axis cs:-0.424,0) rectangle (axis cs:0.424,0.04);
\fill[densityred, opacity=0.35] (axis cs:-1.412,0.045) rectangle (axis cs:1.412,0.085);
\draw[distblue, dashed] (axis cs:-0.424,0.3) -- (axis cs:-0.424,0);
\draw[distblue, dashed] (axis cs:0.424,0.7) -- (axis cs:0.424,0);
\draw[densityred, dashed] (axis cs:-1.412,0.3) -- (axis cs:-1.412,0.045);
\draw[densityred, dashed] (axis cs:1.412,0.7) -- (axis cs:1.412,0.045);
\end{axis}
\end{tikzpicture}
\caption{Why $\SE(\text{median})\propto 1/f(m)$. The fraction of the sample below
any point wobbles by about $\delta=0.5/\sqrt n$ (orange band). Reading that band through
a steep CDF gives a narrow range of medians (blue); through a flat CDF a wide
one (red). The CDF's slope at $m$ is the density $f(m)$.}
\label{fig:sd-cdfmed}
\end{figure}

\begin{example}[How much data does the median cost?]
Normal data, $\sigma=20$, $n=400$: $\SE(\Xbar)=20/20=1.00$, while
$\SE(\text{median})\approx 1.2533$. Since SE scales as $1/\sqrt n$, the median needs
$1.2533^2=\pi/2\approx 1.571$ times as much data to match the mean: about 628
observations instead of 400.
\end{example}

That is the price for Normal data. For skewed or outlier-prone data the
comparison can flip. In our latency population, $m\approx 40.8$ ms and
$f(m)=\tfrac{1}{30}e^{-\ln 2}=1/60$, so $\SE(\text{median})\approx 60/(2\sqrt n)=30/\sqrt n$,
exactly the same as the mean's $\sigma/\sqrt n$. With a single wild outlier
($120$ ms in the eight latencies above) the median barely moves while the mean
and $s$ both jump. Robustness is often worth the efficiency cost.

\begin{trapbox}
The median formula needs $f(m)$, which you do not know and which is hard to
estimate from a small sample. For discrete or heavily tied data (ratings 1 to 5,
latencies rounded to 10 ms) the formula can be badly off. This is exactly the
situation the bootstrap was invented for.
\end{trapbox}

\subsection{Most statistics have no tidy formula}

The mean, a proportion, and a difference of means have exact, simple SEs. The
median has only an approximate one that depends on an unknown density. For the
p99 latency, a trimmed mean, a correlation, a ratio of two means, or the AUC of a
classifier (the area under its ROC curve, a ranking-quality score between 0.5
and 1), a usable closed form is rare or messy. Chapter 8 gives a single
recipe that works for nearly all of them: approximate the sampling distribution by
resampling your own data (the bootstrap).

\begin{figure}[H]
\centering
\begin{tikzpicture}[
  q/.style={draw, rounded corners, fill=histfill!25, align=center, font=\small, minimum height=0.75cm},
  a/.style={draw, rounded corners, fill=targetgreen!15, align=center, font=\small, minimum height=0.75cm},
  b/.style={draw, rounded corners, fill=bootorange!25, align=center, font=\small, minimum height=0.75cm},
  arr/.style={-{Stealth}, thick}]
\node[q] (s) at (0,0) {Which statistic?};
\node[a] (m) at (-4.5,-1.6) {mean\\$s/\sqrt n$};
\node[a] (p) at (-1.5,-1.6) {proportion\\$\sqrt{\hat p(1-\hat p)/n}$};
\node[a] (d) at (1.7,-1.6) {difference\\$\sqrt{\SE_A^2+\SE_B^2}$};
\node[a] (md) at (4.8,-1.6) {median\\$1/(2\hat f(m)\sqrt n)$};
\node[b] (o) at (0,-3.3) {anything else (p99, ratio, correlation, AUC)\\or a formula you do not trust: bootstrap (Chapter 8)};
\foreach \t in {m,p,d,md} \draw[arr] (s) -- (\t);
\draw[arr] (s.south) -- (o.north);
\draw[arr, dashed, gray] (md) |- (o.east);
\end{tikzpicture}
\caption{Choosing a standard error. Closed forms exist for a handful of
statistics; everything else, and any case where the formula's assumptions look
shaky, goes to the bootstrap.}
\label{fig:sd-choose}
\end{figure}

\section{Finite population correction}
\label{sec:sd-fpc}

All formulas so far assume independent draws: sampling with replacement, or from a
population so large that it makes no difference. When you sample \emph{without}
replacement a large fraction $n/N$ of a finite population of size $N$, each draw
removes a unit and the sample covers more of the population than independent draws
would. The SE shrinks by the \textbf{finite population correction}:
\[
\SE(\Xbar)=\frac{\sigma}{\sqrt n}\sqrt{\frac{N-n}{N-1}}\approx\frac{\sigma}{\sqrt n}\sqrt{1-\frac nN}.
\]

\begin{example}[Surveying a company]
A company has $N=2000$ employees; you survey $n=500$ chosen at random without
replacement. The correction is $\sqrt{1500/1999}\approx 0.866$, so the SE is
13\% smaller than the textbook formula. Survey all 2000 and the correction is 0:
there is no sampling error left, only measurement error.
\end{example}

\begin{figure}[H]
\centering
\begin{tikzpicture}
\begin{axis}[width=10cm, height=5cm, xmin=0, xmax=1, ymin=0, ymax=1.08,
  xlabel={sampling fraction $n/N$}, ylabel={correction factor}, axis lines=left,
  xtick={0,0.25,0.5,0.75,1}, ytick={0,0.5,0.866,1},
  xticklabel style={/pgf/number format/fixed}]
\addplot[structurepurple, very thick, domain=0:1, samples=150] {sqrt(1-x)};
\fill[targetgreen!25, opacity=0.6] (axis cs:0,0) rectangle (axis cs:0.05,1.08);
\draw[gray, dashed] (axis cs:0.25,0) -- (axis cs:0.25,0.866) -- (axis cs:0,0.866);
\fill[quotientorange] (axis cs:0.25,0.866) circle (2.5pt);
\node[anchor=west, font=\small] at (axis cs:0.27,0.9) {survey 500 of 2000};
\node[anchor=west, font=\small] (ign) at (axis cs:0.32,0.3) {below 5\%: ignore it};
\draw[-{Stealth}, gray] (ign.west) -- (axis cs:0.055,0.3);
\end{axis}
\end{tikzpicture}
\caption{The finite population correction $\sqrt{1-n/N}$. It stays above $0.97$
while you sample less than 5\% of the population (green), which covers almost all
practical cases. It only matters when you sample a large share.}
\label{fig:sd-fpc}
\end{figure}

\begin{connectionbox}
In machine learning, training usually computes the gradient (the direction
that most reduces the loss) on a minibatch, a small random subset of the
training examples. The minibatch gradient is an average of per-example
gradients, so it is a sample mean. Theorem~\ref{thm:sd-xbar}
is why gradient noise has SD proportional to $1/\sqrt{B}$ for batch size $B$, and
why quadrupling the batch only halves the noise. The same $\sigma/\sqrt n$ governs the
error bars on a test-set accuracy: a classifier scored on 1000 examples at 90\%
accuracy has $\SE=\sqrt{0.9\cdot 0.1/1000}\approx 0.0095$, so a 0.5-point
``improvement'' between two models is well inside the noise.
\end{connectionbox}

\section{Summary}
\label{sec:sd-summary}

\begin{center}
\renewcommand{\arraystretch}{1.35}
\begin{tabular}{lll}
\toprule
Statistic & Standard error & Notes\\
\midrule
mean $\Xbar$ & $\sigma/\sqrt n$, estimated by $s/\sqrt n$ & exact for iid data\\
proportion $\hat p$ & $\sqrt{p(1-p)/n}$ & at most $0.5/\sqrt n$\\
difference $\Xbar_A-\Xbar_B$ & $\sqrt{\SE_A^2+\SE_B^2}$ & independent groups only\\
median & $1/(2f(m)\sqrt n)$; Normal: $1.2533\,\sigma/\sqrt n$ & large $n$\\
other (p99, ratios, \dots) & no simple formula & bootstrap, Chapter 8\\
\bottomrule
\end{tabular}
\end{center}

\begin{itemize}
  \item $\E$ is linear; $\Var(aX+b)=a^2\Var(X)$; variances of independent variables
  add (for sums and differences). SDs do not add.
  \item A statistic is a random variable. Its distribution over repeated samples is
  the sampling distribution; its SD is the standard error.
  \item $\E[\Xbar]=\mu$ and $\Var(\Xbar)=\sigma^2/n$. Four times the data halves the
  SE.
  \item SD describes individuals and does not shrink with $n$. SE describes an
  estimate and does.
  \item CLT: $\Xbar$ is approximately $\Normal(\mu,\sigma^2/n)$ for large $n$. Skewness of
  $\Xbar$ is $\gamma/\sqrt n$, so skewed data need larger $n$. For proportions,
  want $np\ge 10$ and $n(1-p)\ge 10$.
  \item LLN: $\Xbar\to\mu$. Early deviations are diluted, not compensated.
  \item Sampling more than about 5\% of a finite population without replacement:
  multiply the SE by $\sqrt{1-n/N}$.
\end{itemize}

Chapter 7 takes the next step: an estimate plus or minus about two standard
errors becomes a confidence interval, with a precise statement of how often such
intervals catch the truth.

\section{Exercises}
\label{sec:sd-exercises}

\begin{enumerate}
  \item $\Var(X)=4$. Find $\Var(3X+2)$ and $\SD(3X+2)$. If $Y$ is independent of
  $X$ with $\Var(Y)=5$, find $\Var(X-Y)$.
  \item A population has $\sigma=12$. What is the SE of the mean for $n=36$? What
  $n$ gives an SE of 1?
  \item A feature is used by 20\% of users. What is $\SE(\hat p)$ from a random
  sample of 100 users? Without knowing $p$, how large must $n$ be to guarantee
  $\SE(\hat p)\le 0.01$?
  \item Two independent estimates have $\SE_A=0.8$ and $\SE_B=0.6$. What is the SE
  of their difference? What answer does the common mistake give?
  \item From the data $4, 8, 6, 5, 7$ compute $\xbar$, $s$ and the estimated SE of
  the mean.
  \item A report says ``mean session length $6.2\pm 0.1$ min (SE), $n=2500$.''
  Estimate the SD of session length. Would a session of 9 minutes be unusual?
  \item For Normal data with $\sigma=20$ and $n=400$, find the SE of the mean and of
  the median. How many observations would the median need to match the mean's
  SE?
  \item A population has skewness 3. Using $\gamma/\sqrt n\le 0.3$ as the criterion,
  how large must $n$ be for the CLT approximation to be adequate?
  \item You survey $n=100$ of the $N=500$ members of a club, without replacement.
  The population SD is 10. Compute the SE of the mean with and without the finite
  population correction.
  \item Explain in two sentences why $\Var(X_1+X_2)\neq\Var(2X_1)$ for two
  independent dice, though both have mean 7.
\end{enumerate}

\subsection*{Answers}

\begin{enumerate}
  \item $\Var(3X+2)=9\cdot 4=36$, $\SD=6$. $\Var(X-Y)=4+5=9$.
  \item $12/\sqrt{36}=2$. Need $12/\sqrt n=1$, so $n=144$.
  \item $\sqrt{0.2\cdot 0.8/100}=0.04$. Worst case $0.5/\sqrt n\le 0.01$ gives
  $n\ge 2500$.
  \item $\sqrt{0.64+0.36}=1.0$. The mistake (adding SEs) gives $1.4$.
  \item $\xbar=6$; deviations $-2,2,0,-1,1$, squares sum to 10; $s=\sqrt{10/4}\approx 1.581$;
  $\widehat{\SE}=1.581/\sqrt 5\approx 0.707$.
  \item $\SD\approx 0.1\sqrt{2500}=5$ min. A 9-minute session is only $(9-6.2)/5\approx 0.56$
  SDs above the mean: perfectly ordinary. (It would be wildly unusual only if you
  wrongly used the SE.)
  \item $\SE(\Xbar)=1$, $\SE(\text{median})\approx 1.2533$. Need $400\cdot\pi/2\approx 628$
  observations.
  \item $3/\sqrt n\le 0.3$ gives $\sqrt n\ge 10$, so $n\ge 100$.
  \item Without: $10/\sqrt{100}=1$. With: $1\cdot\sqrt{400/499}\approx 0.895$.
  \item Independent dice often offset each other, so the sum's variance is
  $35/12+35/12=35/6$. Doubling one die copies its deviation exactly, so the variance
  is $2^2\cdot 35/12=35/3$, twice as large.
\end{enumerate}
